\documentclass[12pt, a4paper, oneside]{book}

\include{preambulo}

\begin{document}




%---------- CAPA -------------





% Não numera a capa
\thispagestyle{empty}

\begin{center}
	% \textsc{\large Universidade Federal do Rio Grande do Norte \\ Centro de Ciências Exata e da Terra \\ Departamento de Matemática \\	Mestrado Profissional em \\ Matemática em Rede Nacional - PROFMAT}\\
%\includegraphics{logo.pdf}\\
\begin{center}
    \includegraphics[width=.2\linewidth]{Figuras/brasao_gradiente.png}
\end{center}


\begin{center}
    UNIVERSIDADE FEDERAL DO RIO GRANDE DO NORTE \\
    CENTRO DE CIÊNCIAS EXATAS E DA TERRA \\
    DEPARTAMENTO DE FÍSICA \\
    BACHARELADO EM FÍSICA
\end{center}

\vspace{2cm}

%-------------Autor------------

\begin{center}
    \large Autor
\end{center}

\vspace{4cm}

%-----------Título----------------

\textbf{\Large{Título}}\\

\vspace{6cm}



\large Natal - RN \\
Janeiro de 2019
\end{center}





%---------- FOLHA DE ROSTO -------------




\newpage
\thispagestyle{empty}

\begin{center}
	\textbf{\large Autor}\\
	\vspace{2.5cm}
	\textbf{\Large{Título}}
\end{center}
	\vspace{2.2cm}
\begin{flushright}
\begin{minipage}{8cm}
Monografia de Graduação apresentada ao Departamento
de Física Teórica e Experimental do Centro de Ciências Exatas e da Terra da Universidade Federal do Rio Grande do Norte como requisito parcial para a obtenção do grau de bacharel em Física.
\end{minipage}
\end{flushright}
 
\vspace{2cm}

\begin{center}
   \large Orientador: \\
    Prof. Dr. Orientador
\end{center}

\vspace{1cm}

\begin{center}
    Universidade Federal do Rio Grande do Norte - UFRN \\
    Departamento de Física - DF
\end{center}

\vspace{2cm}
\begin{center}
\textbf{Natal - RN} \\
Janeiro de 2019
\end{center}




\begin{comment}

%---------- BANCA EXAMINADORA -------------
%\includepdf{Assinaturas/Assinaturas.pdf}
\newpage
\thispagestyle{empty}
	\begin{center}
		\textsc{ Universidade Federal do Rio Grande do Norte \\ Centro de Ciências Exata e da Terra \\ Departamento de Matemática \\	Mestrado Profissional em Matemática em Rede Nacional - PROFMAT}\\
	
	\vspace{3cm}
	\textsc{George Homer Barbosa de Medeiros}\\
	\vspace{3cm}
	\textsc{\Huge{Coloque aqui o título da sua dissertação}}
	\end{center}
	\vspace{4cm}
Comissão Examinadora:\vspace{0,5cm}\\
Prof. Dr. AAAAAAAAAAAAAA (UFRN - Orientador)\\
Prof. Dr. BBBBBBBBBBBBBB  (UFRN - Membro interno)\\
Prof. Dr. CCCCCCCCCCCCCC (UFRN - Membro interno)\\
Prof. Dr. DDDDDDDDDDDDDD  (IFCE - Membro externo)\\
\vspace{2cm}

\begin{center}
	\textbf{Natal - RN} \\
	Janeiro de 2019
\end{center}

\end{comment}



\frontmatter

\pagenumbering{roman}


%------------Agradecimentos------------

\include{agradecimentos}

%------------Citação-------------------
%\addcontentsline{toc}{chapter}{Foreword}
\include{citacao}


%--------------Resumo-------------------
\addcontentsline{toc}{chapter}{Resumo}

\pagestyle{myheadings} % tira texto do cabeçalho
\newpage

\chapter*{Resumo}
\thispagestyle{myheadings}
\noindent
Digite aqui seu resumo dfd fdf dfdfdg d hfgh gçhkg hjfhj fihjf gihj gihj fgihjijh dfd fdf dfdfdg d hfgh gçhkg hjfhj fihjf gihj gihj fgihjijh dfd fdf dfdfdg d hfgh gçhkg hjfhj fihjf gihj gihj fgihjijh dfd fdf dfdfdg d hfgh gçhkg hjfhj fihjf gihj gihj fgihjijh dfd fdf dfdfdg d hfgh gçhkg hjfhj fihjf gihj gihj fgihjijh dfd fdf dfdfdg d hfgh gçhkg hjfhj fihjf gihj gihj fgihjijh dfd fdf dfdfdg d hfgh gçhkg hjfhj fihjf gihj gihj fgihjijh.
 
\noindent  {\bf Palavras-chave:} Aaaa, Bbbb, Cccc, Dddd.


%-------------Abstract------------------

\newpage
\addcontentsline{toc}{chapter}{Abstract}

\chapter*{Abstract}
\thispagestyle{myheadings}
\noindent
Digite aqui seu resumo dfd fdf dfdfdg d hfgh gçhkg hjfhj fihjf gihj gihj fgihjijh dfd fdf dfdfdg d hfgh gçhkg hjfhj fihjf gihj gihj fgihjijh dfd fdf dfdfdg d hfgh gçhkg hjfhj fihjf gihj gihj fgihjijh dfd fdf dfdfdg d hfgh gçhkg hjfhj fihjf gihj gihj fgihjijh dfd fdf dfdfdg d hfgh gçhkg hjfhj fihjf gihj gihj fgihjijh dfd fdf dfdfdg d hfgh gçhkg hjfhj fihjf gihj gihj fgihjijh dfd fdf dfdfdg d hfgh gçhkg hjfhj fihjf gihj gihj fgihjijh.

\noindent  {\bf Keywords:} Aaaa, Bbbb, Cccc, Dddd.



%-------------Lista de figuras----------
\newpage
\addcontentsline{toc}{chapter}{Lista de Figuras}
\listoffigures
\thispagestyle{myheadings}


%-------------Lista de Tabelas----------
\newpage
\addcontentsline{toc}{chapter}{Lista de Tabelas}
\listoftables
\thispagestyle{headings}

%\mainmatter

%-------------Índice--------------------


\newpage


\tableofcontents
\thispagestyle{myheadings}

%-------------Introdução----------------
\mainmatter


\newpage

\pagenumbering{arabic}
\pagestyle{myheadings}

%-------------Capítulo 1-----------------
\include{capitulo1}

%-------------Capítulo 2-----------------
\include{capitulo2}

%-------------Conclusão-----------------
\include{conclusao}

%\nocite{khan,da2005aplicaccao,gama2004,lior2014,dados1}

%-------------Bibliografia------------------
\newpage
\begin{thebibliography}{99}
	
	\bibitem{AN} ANDREESCU, T., ANDRICA, D., {\bf 360 Problems for Mathematical Contests}, GIL Publishing House, Zalău, 2003.
	
	\bibitem{JE}  ANDREESCU T., CIRTOJAN,E V., Dospinescu G., Lascu M., {\bf Old and New Inequalities}, GIL Publishing House, 2004.
	
	\bibitem{VO} CAN, V.Q.B,Pohoat$\breve{a}$, C., {\bf Old and New inequalities, Volume 2,} GIL Publishing House, 2008.
	
	\bibitem{CV} CVETKOVSKI, Zdravko; {\bf Inequalities - Theorems, Techniques and Selected Problems.}	editora Springer-Verlag Berlin Heidelberg, ano 2012.
	
	\bibitem{Bartle} BARTLE, Robert G.; {\bf Introduction to Real Analysis.}
	editora Willey, ano 2011.
	
\end{thebibliography}


%-----------APENDICES---------------------
\renewcommand{\refname}{APÊNDICES}
\addcontentsline{toc}{chapter}{APÊNDICES}
\appendix % Cue to tell LaTeX that the following "chapters" are Appendices

\include{Apendices/apendiceA}
\include{Apendices/apendiceB}

\end{document}